\documentclass[a3paper,landscape,10pt]{article}
\usepackage[margin=10mm]{geometry}
\usepackage[french]{babel}
\usepackage{amsmath,amssymb}
\makeatletter
\def\input@path{{../../../Style/}}
\makeatother
\NeedsTeXFormat{LaTeX2e}
\ProvidesPackage{TLPosterCarteMentale}[2026/08/03 Posters et cartes mentales SI]

\RequirePackage{tikz}
\RequirePackage[most]{tcolorbox}
\RequirePackage{xcolor}
\RequirePackage{amsmath,amssymb}
\RequirePackage{graphicx}
\usetikzlibrary{positioning,arrows.meta,calc,shapes.geometric,mindmap,backgrounds,fit}

\definecolor{TLPosterBlue}{RGB}{24,86,141}
\definecolor{TLPosterLight}{RGB}{234,243,250}
\definecolor{TLPosterAccent}{RGB}{232,126,34}
\definecolor{TLPosterGreen}{RGB}{52,152,102}
\definecolor{TLPosterRed}{RGB}{192,57,43}
\definecolor{TLPosterGray}{RGB}{80,90,100}

\newcommand{\TLPosterTitre}[2]{%
  \begin{tcolorbox}[enhanced,boxrule=0pt,arc=0pt,colback=TLPosterBlue,left=6mm,right=6mm,top=4mm,bottom=4mm]
    {\color{white}\bfseries\LARGE #1}\hfill{\color{white}\large #2}
  \end{tcolorbox}%
}

\newtcolorbox{TLPosterBloc}[2][]{
  enhanced,breakable,
  colback=TLPosterLight,
  colframe=TLPosterBlue,
  colbacktitle=TLPosterBlue,
  coltitle=white,
  fonttitle=\bfseries,
  title={#2},
  boxrule=0.8pt,
  arc=2mm,
  left=3mm,right=3mm,top=2mm,bottom=2mm,
  #1
}

\newtcolorbox{TLPosterFormule}[1][]{
  enhanced,
  colback=white,
  colframe=TLPosterAccent,
  boxrule=1pt,
  arc=2mm,
  fontupper=\bfseries,
  halign=center,
  #1
}

\newcommand{\TLPosterMethode}[1]{%
  \begin{tcolorbox}[enhanced,colback=TLPosterGreen!8,colframe=TLPosterGreen,title=Méthode,fonttitle=\bfseries]
  #1
  \end{tcolorbox}%
}

\newcommand{\TLPosterAttention}[1]{%
  \begin{tcolorbox}[enhanced,colback=TLPosterRed!6,colframe=TLPosterRed,title=Point de vigilance,fonttitle=\bfseries]
  #1
  \end{tcolorbox}%
}

\tikzset{
  TLmindroot/.style={concept,concept color=TLPosterBlue,text=white,font=\bfseries\large,minimum size=34mm},
  TLmindlevel1/.style={level 1 concept/.append style={font=\bfseries,minimum size=23mm,sibling angle=45}},
  TLmindlevel2/.style={level 2 concept/.append style={font=\small,minimum size=17mm}},
  TLarrow/.style={-{Stealth[length=3mm]},thick,draw=TLPosterBlue},
  TLbox/.style={rounded corners=2mm,draw=TLPosterBlue,fill=TLPosterLight,align=center,inner sep=3mm}
}

\newenvironment{TLCarteMentale}[1]{%
  \begin{tikzpicture}[mindmap,grow cyclic,concept color=TLPosterBlue,every node/.style={concept},TLmindlevel1,TLmindlevel2]
  \node[TLmindroot]{#1}
}{;
  \end{tikzpicture}
}

\newcommand{\TLBranche}[3][]{child[concept color=#2]{node[#1]{#3}}}

\endinput

\IfFileExists{TLRDM.sty}{\ProvidesPackage{TLRDM}[2026/08/03 Outils graphiques RDM]
\RequirePackage{tikz}
\usetikzlibrary{arrows.meta,calc,patterns,decorations.pathreplacing}

\newcommand{\TLPoutreReperee}{%
\begin{center}\begin{tikzpicture}[>=Latex,scale=.9]
\draw[very thick] (0,0)--(8,0);
\draw[->] (0,-.8)--(1.4,-.8) node[right]{$\vec{x}$};
\draw[->] (0,-.8)--(0,.6) node[above]{$\vec{y}$};
\fill (0,0) circle (1.5pt) node[above left]{$A$};
\fill (8,0) circle (1.5pt) node[above right]{$B$};
\draw[dashed] (4,-.65)--(4,.65);
\node[above] at (4,.65) {section droite $S$};
\end{tikzpicture}\end{center}}

\newcommand{\TLCoupurePoutre}{%
\begin{center}\begin{tikzpicture}[>=Latex,scale=.9]
\draw[very thick] (0,0)--(7,0);
\draw[dashed] (3.5,-.8)--(3.5,.8);
\draw[->,thick] (3.5,0)--(4.8,0) node[right]{$N\vec{x}$};
\draw[->,thick] (3.5,0)--(3.5,1.3) node[above]{$T_y\vec{y}$};
\draw[->,thick] (3.1,.45) arc[start angle=145,end angle=-145,radius=.45];
\node at (2.5,.9) {$M_z$};
\node[below] at (3.5,-.85) {coupure en $G$};
\end{tikzpicture}\end{center}}

\newcommand{\TLDiagrammeTraction}{%
\begin{center}\begin{tikzpicture}[>=Latex]
\draw[fill=gray!15] (0,-.35) rectangle (6,.35);
\draw[->,very thick] (0,0)--(-1.3,0) node[left]{$F$};
\draw[->,very thick] (6,0)--(7.3,0) node[right]{$F$};
\draw[decorate,decoration={brace,mirror,amplitude=5pt}] (0,-.65)--(6,-.65) node[midway,below=5pt]{$L$};
\end{tikzpicture}\end{center}}

\newcommand{\TLDiagrammeFlexion}{%
\begin{center}\begin{tikzpicture}[>=Latex,scale=.9]
\draw[very thick] (0,0)--(7,0);
\draw (0,0)--(-.35,-.55)--(.35,-.55)--cycle;
\draw (7,0)--(6.65,-.55)--(7.35,-.55)--cycle;
\draw[->,very thick] (3.5,1.4)--(3.5,.05) node[midway,right]{$F$};
\draw[dashed] (0,0) .. controls (2.2,-.25) and (4.8,-.25) .. (7,0);
\node[below] at (3.5,-.45) {déformée};
\end{tikzpicture}\end{center}}

\newcommand{\TLDistributionContrainteFlexion}{%
\begin{center}\begin{tikzpicture}[>=Latex,scale=.9]
\draw[fill=gray!10] (-1,-1.6) rectangle (1,1.6);
\draw[dashed] (-1,0)--(2.8,0) node[right]{fibre neutre};
\draw[thick] (1,-1.6)--(2.5,0)--(1,1.6);
\node[right] at (2.5,.25) {$\sigma=0$};
\node[right] at (1,1.35) {compression};
\node[right] at (1,-1.35) {traction};
\end{tikzpicture}\end{center}}
}{}
\pagestyle{empty}
\begin{document}
\TLPosterTitre{Résistance des matériaux}{Poutres, contraintes, déformations et dimensionnement}
\begin{multicols}{3}

\begin{TLPosterBloc}{Objectifs}
Vérifier qu'une pièce résiste aux sollicitations, que ses déformations restent admissibles et que sa stabilité est assurée.
\end{TLPosterBloc}

\begin{TLPosterBloc}{Modèle poutre}
Une poutre possède une longueur grande devant les dimensions de sa section. Sa ligne moyenne est décrite par une abscisse $x$ et une section droite $S$.
\end{TLPosterBloc}

\begin{TLPosterBloc}{Hypothèses usuelles}
Matériau homogène et isotrope, petites perturbations, principe de Saint-Venant et hypothèse de Navier--Bernoulli pour la flexion des poutres élancées.
\end{TLPosterBloc}

\begin{TLPosterBloc}{Torseur de cohésion}
Une coupure en une section fait apparaître les efforts intérieurs : effort normal $N$, efforts tranchants $T_y,T_z$, moment de torsion $M_t$ et moments fléchissants $M_{fy},M_{fz}$.
\end{TLPosterBloc}

\begin{TLPosterBloc}{Sollicitations simples}
Traction-compression : $N$. Cisaillement : $T$. Torsion : $M_t$. Flexion : $M_f$. Une sollicitation composée associe plusieurs composantes.
\end{TLPosterBloc}

\begin{TLPosterBloc}{Contrainte et déformation}
La contrainte traduit l'intensité des efforts internes. La déformation normale vaut $\varepsilon=\Delta L/L$ et la distorsion est notée $\gamma$.
\end{TLPosterBloc}

\begin{TLPosterBloc}{Loi de Hooke}
Dans le domaine élastique linéaire :
\[
\sigma=E\varepsilon,\qquad \tau=G\gamma,
\]
où $E$ est le module de Young et $G$ le module de cisaillement.
\end{TLPosterBloc}

\begin{TLPosterBloc}{Traction-compression}
Pour une section sollicitée uniformément :
\[
\sigma=\frac{N}{S},\qquad \Delta L=\frac{NL}{ES}.
\]
Le signe de $N$ distingue traction et compression.
\end{TLPosterBloc}

\begin{TLPosterBloc}{Cisaillement simple}
Une estimation moyenne est donnée par
\[
\tau_{\mathrm{moy}}=\frac{T}{S}.
\]
La répartition réelle dépend de la géométrie de la section.
\end{TLPosterBloc}

\begin{TLPosterBloc}{Torsion d'une section circulaire}
Pour un arbre circulaire :
\[
\tau(r)=\frac{M_t r}{J},\qquad
\theta=\frac{M_t L}{GJ},
\]
où $J$ est le moment quadratique polaire.
\end{TLPosterBloc}

\begin{TLPosterBloc}{Flexion simple}
La contrainte normale varie linéairement dans la section :
\[
\sigma=-\frac{M_f y}{I}.
\]
Elle est maximale sur les fibres les plus éloignées de la fibre neutre.
\end{TLPosterBloc}

\begin{TLPosterBloc}{Déformée en flexion}
Pour de petites déformations :
\[
EI\,y''(x)=M_f(x).
\]
Deux intégrations, complétées par les conditions aux limites, donnent rotation et flèche.
\end{TLPosterBloc}

\begin{TLPosterBloc}{Moments quadratiques}
Le moment quadratique $I$ dépend de la géométrie et de l'axe considéré. Le théorème de Huygens permet de transporter un axe : $I=I_G+Sd^2$.
\end{TLPosterBloc}

\begin{TLPosterBloc}{Critères de résistance}
En sollicitation simple, vérifier $|\sigma_{\max}|\leq \sigma_{\mathrm{adm}}$ ou $|\tau_{\max}|\leq \tau_{\mathrm{adm}}$. En état multiaxial, utiliser un critère adapté comme Tresca ou von Mises.
\end{TLPosterBloc}

\begin{TLPosterBloc}{Rigidité et sécurité}
Le dimensionnement doit satisfaire à la fois la résistance, la rigidité et un coefficient de sécurité tenant compte des incertitudes, dispersions et conditions d'usage.
\end{TLPosterBloc}

\begin{TLPosterBloc}{Démarche d'étude}
Modéliser la poutre, déterminer les actions extérieures, calculer les réactions, établir le torseur de cohésion, tracer les diagrammes, calculer contraintes et déformations, puis conclure sur les critères.
\end{TLPosterBloc}

\end{multicols}
\end{document}
